\textit{Segmentasi instansi inti sel pada histopatologi H\&E merupakan hal yang esensial untuk tugas patologi komputasional seperti penilaian tumor-infiltrating lymphocyte (TILs) dan grading kanker. Namun, inti sel yang saling tumpang tindih dan ketergantungan metode sebelumnya pada pemrosesan pasca-ekstraksi buatan (NMS, watershed) menimbulkan tantangan geometris dan komputasional yang menghambat penerapan pada perangkat edge dengan sumber daya terbatas. Metode yang ada saat ini masih berbasis mask dan bergantung pada pemrosesan pasca-ekstraksi tersebut untuk memisahkan instansi yang saling bersentuhan, sehingga tidak cocok untuk inferensi edge yang ringan.}

\textit{Penelitian ini memperkenalkan RayCastED, model segmentasi instansi berbasis kontur yang ringan, sepenuhnya konvolusional, dan end-to-end, yang memprediksi poligon star-convex langsung dari grid anchor tanpa pemrosesan pasca-ekstraksi apa pun. RayCastED mengadopsi backbone deteksi turunan YOLO26 dengan dual-label assignment untuk inferensi tanpa NMS, inter-scale competition untuk menekan deteksi duplikat antar-skala, ResoConv berbasis wavelet untuk pelestarian batas, dan split classification head untuk mengatasi ketidakseimbangan foreground-background. Loss jarak ray analitis yang baru menghitung ulang ray ground-truth dari centroid yang diprediksi, menggabungkan akurasi centroid dengan supervisi batas dan memberikan gradien yang invarian terhadap skala.}

\textit{Evaluasi pada PanNuke dengan 3-fold cross-validation menunjukkan bahwa RayCastED mencapai panoptic quality yang kompetitif dibandingkan baseline yang jauh lebih berat sambil mencapai performa per-jaringan yang paling konsisten, dan mendemonstrasikan inferensi real-time yang layak pada platform embedded Jetson Orin Nano. Hasil ini menetapkan RayCastED sebagai model berbasis kontur tanpa pemrosesan pasca-ekstraksi yang menyatukan akurasi batas, inferensi ringan, dan keterdeployabilitas edge untuk analisis inti sel. Kode dan model tersedia di \url{https://github.com/DanishRitonga/RayCastED}.}

\noindent{Kata kunci} : Segmentasi Instansi Inti Sel, Patologi Komputasional, Poligon Star-convex, Arsitektur Ringan, Deployment Edge
